%% kapa: exam=1 round=28 year=2017 period=1 form=A subject=civil_law subject_ko=민법 seq=33 number=33
%% source: Q-Net 감정평가사 기출문제 (cst00309) artlSeq=5199603 file=extracted/1008/5199603/2017년 제28회 감정평가사 1차시험 문제지/1교시 A형.pdf
%% conversion: kordoc 4.24.51 -> tools/gen1.py (auto); review=proofread-sonnet; flags=none
\begin{question}{33}
관습상 법정지상권에 관한 설명으로 옳지 않은 것은? (다툼이 있으면 판례에 따름)
\begin{choices}
  \item 토지공유자 중 1인이 다른 공유자의 지분 과반수의 동의를 얻어 공유토지 위에 건물을 건축한 후 토지와 건물의 소유자가 달라진 경우, 관습상 법정지상권은 성립하지 않는다.
  \item 강제경매에 있어 관습상 법정지상권이 인정되기 위해서는 매각대금 완납 시를 기준으로 해서 토지와 그 지상 건물이 동일인의 소유에 속하여야 한다.
  \item 관습상 법정지상권자는 토지소유자로부터 토지를 양수한 자에 대하여 등기 없이도 자신의 권리를 주장할 수 있다.
  \item 대지와 건물의 소유자가 건물만을 양도하면서 양수인과 대지에 관하여 임대차계약을 체결한 경우, 특별한 사정이 없는 한 그 양수인은 관습상 법정지상권을 포기한 것으로 본다.
  \item 구분소유적 공유관계에 있는 자가 자신의 특정 소유가 아닌 부분에 건물을 신축한 경우, 관습상 법정지상권이 성립하지 않는다.
\end{choices}
\end{question}
\answer{2}
